\title{LongRoPE: Extending LLM Context Window Beyond 2 Million Tokens}

\begin{abstract}
An extensive context window is a highly sought-after capability in large language models (LLMs). Yet, several challenges—such as the expensive nature of fine-tuning, the limited availability of sufficiently long textual data, and the emergence of detrimental values at novel token positions—have thus far restricted extended context windows to approximately 128k tokens.

This work presents LongRoPE, marking the first approach to broaden the context window of pre-trained LLMs to a substantial \textbf{2048k} tokens, employing at most 1k fine-tuning iterations confined to training lengths up to 256k, all while sustaining the performance level of the original short context window. The proposed method is based on three central innovations: (i) we uncover and leverage two kinds of non-uniformities in positional interpolation via an efficient search procedure, yielding superior fine-tuning initialization and supporting an 8$\times$ extension for inference even without fine-tuning; (ii) we propose a staged extension scheme, which starts by fine-tuning the model at a 256k length, followed by a secondary positional interpolation on this already extended model to ultimately reach a 2048k context window; (iii) we re-calibrate LongRoPE on sequences of length 8k in order to recapture the original model’s short-range capabilities. Comprehensive evaluations on LLaMA2 and Mistral, spanning multiple benchmarks, validate the utility of our approach. Models that undergo extension with LongRoPE require only slight modifications to positional embeddings, keeping the remainder of the original architecture intact and allowing for compatibility with almost all pre-existing optimizations. Code will be released at \texttt{https://github.com/microsoft/LongRoPE}

\section{Introduction}

Large Language Models (LLMs) have achieved impressive results across a range of tasks~\cite{gpt4,llama2}, yet they remain constrained by a relatively short context window—such as the 4096-token limitation found in LLaMA2~\cite{llama2}. When inputs exceed this window, the performance of LLMs deteriorates, primarily because the models encounter positional inputs beyond those present during training. This issue creates difficulties for critical applications, including in-context learning with large numbers of examples~\cite{Huang2023FewerIM} and for the effective operation of LLM agents~\cite{park2023generative,madaan2023selfrefine}.

Recent studies have demonstrated that it is possible to enlarge the context window of a pre-trained LLM to approximately 128k tokens by fine-tuning with longer sequences~\cite{longlora,pi,yarn,zhang2024soaring,liu2023scaling}. However, scaling the context window further presents several key challenges. First, incorporating previously unseen positional indices results in significant distributional shifts, often introducing harmful values and causing issues that hinder the convergence of fine-tuning~\cite{pi}. This complication intensifies in scenarios where the window is expanded dramatically (for instance, from 4k to over $>$1000k), as this process brings in over 90\% new positions. Second, effective fine-tuning typically necessitates training samples whose lengths match the extended context window. Unfortunately, current corpora contain only a small number of ultra-long texts, particularly those above 1000k tokens. In addition, the cost associated with training on such extensive sequences is very high, demanding vast GPU resources and extensive training times. Third, once the context window is extended to extreme lengths, the model’s attention mechanism may become diluted, distributing focus over a much broader span of tokens and, as a consequence, impairing performance on shorter input contexts~\cite{pi}.

A common strategy for addressing the first challenge involves using interpolated RoPE positional embeddings~\cite{rope,pi}, which adjust novel position indices to fit within the pretrained range, as illustrated in Fig.\ref{fig:overview}. The Position Interpolation (PI) method~\cite{pi} achieves this by linearly scaling RoPE's rotary angles according to the expansion factor. NTK~\cite{ntk,dynamicntk}, by contrast, proposes non-uniform interpolation and extrapolation techniques along different RoPE dimensions. Meanwhile, YaRN~\cite{yarn} divides RoPE's dimensions based on frequency into three categories, applying distinct treatments—extrapolation, NTK interpolation, and linear interpolation—to each group. Nonetheless, the positional embedding structure within Transformers demonstrates a \emph{complex, heterogeneous distribution of information entropy}. Existing solutions tend to overlook this intricate variability, failing to exploit it and resulting in diminished information preservation, which in turn restricts achievable context window lengths.

Section~\ref{sec:analysis} empirically demonstrates two principal observations: \textbf{(1)} Optimal positional interpolation must address two sources of non-uniformity—differences in RoPE dimensions and in token positions. Less interpolation is preferable for lower RoPE dimensions and tokens appearing early in the sequence; however, the best configuration is contingent upon the specific extended length being targeted. \textbf{(2)} By integrating these non-uniformities into the design of positional interpolation, it becomes possible to preserve vital information inherent to the original RoPE, particularly in the crucial dimensions and initial token positions. This approach effectively reduces information loss resulting from interpolation and therefore establishes a stronger initialization for subsequent fine-tuning. Additionally, in settings without fine-tuning, this technique enables an extension of up to 8$\times$ the original sequence length.

Building upon these insights, we present \textbf{LongRoPE}, a robust approach designed to expand the LLM context window past 2 \emph{million} tokens. The design of LongRoPE hinges on three principal advancements. To start, {LongRoPE} capitalizes on multidimensional non-uniformities inherent in positional interpolation, determining optimal rescale parameters for RoPE rotation angles in each RoPE dimension as dictated by token placement. Since the complexity of identifying these rescale factors grows exponentially with the intended window size, {LongRoPE} employs an evolutionary search algorithm supplemented by two optimization strategies that significantly enhance the efficiency of this search. An illustration of the discovered rescaled RoPE is provided in Fig.~\ref{fig:overview}.

Subsequently, {LongRoPE} adopts a resource-efficient, staged extension approach that enables scaling the context window to 2048k tokens—accomplished without directly fine-tuning on ultra-long texts, which are limited in availability. The process starts by identifying a viable context window of 256k within the pre-trained LLM, followed by targeted fine-tuning at this length. Owing to our non-uniform positional interpolation technique, which supports up to an 8$\times$ increase in context length without requiring additional fine-tuning, a second search is performed to determine appropriate RoPE rescale factors for the already fine-tuned, extended model. Through this method, we are able to achieve a 2048k context window for both LLaMA2 and Mistral~\cite{mistral}.

To address potential declines in performance for shorter context windows, {LongRoPE} further modifies the RoPE rescale factors within the extended LLM. Following an analogous approach used when scaling up from 256k to 2048k, we instead reduce the context window to 4k and 8k on the LLM that has been fine-tuned at 256k, employing our search algorithm to minimize the extent of positional interpolation. During inference, when the sequence length falls below 8k, RoPE is adjusted using the rescale factors identified through this search process.

A comprehensive set of experiments involving multiple LLMs and a range of long-context evaluation tasks validate the efficacy of our approach. Our results indicate that LongRoPE consistently preserves low perplexity across evaluation lengths from 4k up to 2048k, attains more than 90\% accuracy in passkey retrieval, and achieves performance on par with state-of-the-art on established benchmarks constructed within a 4096-token context window. Furthermore, LongRoPE is compatible with any LLM utilizing RoPE embeddings. The codebase and LongRoPE-2048k model checkpoints will be made publicly available.

\section{Non-uniformity in Positional Interpolation}
\label{sec:analysis}

\subsection{Preliminary}

Transformer models require explicit positional information, often in the form of position embedding, to represent the order of input tokens. Our work focuses on the RoPE~\cite{rope} position embedding, which is widely used in recent  LLMs. For a token at position index $n$, its corresponding RoPE encoding can be simplified as follows:

\begin{equation}
		\fontsize{7.8}{7.8} \selectfont
	\label{eq:rope}
	\left[ cos(n\theta_0), sin(n\theta_0), cos(n\theta_1), \cdot\cdot\cdot, cos(n\theta_{d/2-1}), sin(n\theta_{d/2-1}) \right]   
\end{equation}

Here, $d$ refers to the embedding dimension, $n\theta_i$ denotes the rotational angle assigned to the token located at position $n$, and $\theta_i=\theta^{-2i/d}$ encodes the rotation frequency. Typically, in RoPE, the parameter $\theta$ is set to 10000.

\textbf{\emph{Context window extension ratio $s$} and positional interpolation}. The ratio $s$ quantifies how much the context window is extended, and is defined as $s=\frac{L'}{L}$, where $L'$ stands for the new, extended context length, and $L$ is the initial context length.

To extend the context window from $L$ to $L'$, current positional interpolation methods suggest downscaling rotation frequencies $\theta_i$ based on extension ratio $s$. Let $\beta=\theta^{2/d}$, and $\lambda$ denote the actual rescale factor related to $s$, we   unify these positional interpolation methods as follows:

\begin{equation}	
	\fontsize{7.5}{7.5} \selectfont
	\label{eq:unified_rope}
	\begin{aligned}
		\left[cos\left(\frac{n}{\lambda(\beta)^0}\right), sin \left(\frac{n}{\lambda(\beta)^0}\right), cos\left(\frac{n}{\lambda(\beta)^1}\right), \cdot\cdot\cdot,  sin \left(\frac{n}{\lambda(\beta)^{d/2-1}}\right) \right]   
	\end{aligned}
\end{equation}

\textbf{Linear positional interpolation (PI)}. The PI~\cite{pi} method proposes linearly interpolating position indices, constrained to the maximum sequence length seen during pre-training. When applying an extension ratio $s$, the corresponding rotation angles for every position are uniformly scaled down by a factor of $\lambda=s$ over each RoPE dimension. Despite this, such compression of position information results in highly compacted representations, making it difficult for the model to differentiate tokens that are near each other. As a consequence, PI typically shows suboptimal performance when used with larger extension ratios.

\textbf{NTK-based interpolation and extrapolation}. ~\cite{ntk,dynamicntk} examine RoPE from the viewpoint of information encoding and leverage Neural Tangent Kernel (NTK) theory~\cite{jacot2018neural,tancik2020fourier}. In order to address the issue of crowded positions in PI, their approach involves redistributing the interpolation burden across the various dimensions of RoPE. Specifically, lower (corresponding to higher frequency) dimensions undergo less scaling, while higher (corresponding to lower frequency) dimensions are scaled more significantly. This adjustment enables both interpolation and extrapolation of positional encodings, with $\lambda=s^{i}$. The enhanced dynamic NTK~\cite{dynamicntk} further modifies the extension ratio adaptively at each position depending on the sequence's current length. In contrast to PI, which requires additional fine-tuning, NTK-based approaches support expansion of context windows even in the absence of fine-tuning; however, such methods generally allow up to only a 4$\times$ increase in context window size.

\textbf{YaRN}~\cite{yarn} offers a notable advancement in the effectiveness of positional interpolation by categorizing RoPE dimensions into three distinct groups based on their frequency characteristics, with each group following a separate interpolation approach. For dimensions corresponding to high frequencies, extrapolation ($\lambda$=1) is used; for those associated with low frequencies, linear interpolation (PI) is applied. Dimensions situated between these extremes are processed with the NTK method. The principal contribution of YaRN is its strategy of grouping RoPE dimensions, though this currently relies on empirical decisions made by humans, which could limit its effectiveness when adapted to novel LLM architectures.

\subsection{Study on Non-uniform Positional Interpolation}

Drawing on insights from NTK and YaRN, we observe that their improvements are attributed to the incorporation of non-linearities, particularly through the treatment of varying frequencies across RoPE dimensions to enable targeted interpolation and extrapolation. Despite these advances, present non-linear strategies are predominantly based on heuristics established by researchers. This prompts two central questions: (1) Does the prevailing approach to positional interpolation represent an optimal solution? (2) Are there additional, as-yet-unstudied forms of non-linearity?

In order to investigate these questions, we employ evolution search (refer to Sec.~\ref{sec:method}) to identify more effective non-uniform positional interpolation strategies for LLaMA2-7B. This process is directed by perplexity measurements on 5 randomly chosen instances from the PG19~\cite{pg19} validation dataset. Our experimental evaluation leads to several notable insights.

\textbf{Finding 1}: \textit{RoPE dimensions exhibit substantial non-uniformities, which are not effectively handled by current positional interpolation methods.}

We determine the optimal value of $\lambda$ for each RoPE dimension as described in Eq.~\ref{eq:unified_rope}. Table~\ref{tbl:exp1} presents a comparison of perplexity scores for LLaMA2-7B across various approaches, evaluated on the PG19 and Proof-pile~\cite{proof-pile} test datasets, all without any fine-tuning. Our method of searching for solutions demonstrates marked improvements over existing techniques, implying that previously used linear (PI) and non-uniform (Dynamic-NTK and YaRN) interpolation strategies do not yield optimal results. Interestingly, on the PG19 dataset, YaRN exhibits poorer performance relative to PI and NTK, which can be attributed to its failure to achieve the intended context window size in non-fine-tuned language models. Specifically, for an 8k context window, the perplexity with YaRN increases abruptly after reaching 7k.

In our investigation, the rescaling factors $\lambda$ in Eq.~\ref{eq:unified_rope} are not uniform; instead, they vary, unlike the constant scale $s$ used in PI, the formulaic approach of NTK, or YaRN's group-wise method. The introduction of these heterogeneous scaling factors leads to notable gains in the language modeling capabilities of LLaMA2—specifically, in terms of lower perplexity for 8k and 16k context lengths—achieved without additional fine-tuning. This performance boost can be attributed to the enhanced ability of the modified positional embedding to maintain the structure of the original RoPE, particularly along the critical dimensions, thereby minimizing the challenge faced by LLMs in distinguishing among nearby token positions.

\textbf{Finding 2}: \textit{RoPE for the initial tokens in the input sequence should be extrapolated with less interpolation.} 

We propose that for the first $\hat{n}$ tokens in the input sequences, the RoPE mechanism should minimize interpolation. This suggestion stems from the fact that these early tokens tend to accumulate higher attention scores and thus play a vital role within the attention modules, a phenomenon previously documented in Streaming LLM~\cite{streamingllm} and LM-Infinite~\cite{han2023lminfinite}. To investigate this claim, we increased the context window to 8k and 16k tokens via PI and NTK, but left the initial $\hat{n}$ tokens (with values 0, 2, ..., 256) unaltered by interpolation. The setting with $\hat{n}=0$ corresponds to the standard PI and NTK configurations. As presented in Table~\ref{tbl:tokenposition}, two central findings emerge: \textbf{(1)} Exempting the initial tokens from position interpolation consistently enhances performance; \textbf{(2)} The ideal value for $\hat{n}$, denoting the number of early tokens excluded from interpolation, varies depending on the intended length of context window extension.

\textbf{Finding 3}: \textit{Non-uniform positional interpolation effectively extends LLM context window in both fine-tuning and non-fine-tuning settings.} 

Although our non-uniform position interpolation strategy already yields considerable gains in extension performance at 8k and 16k context lengths without additional fine-tuning, achieving robust results for even longer contexts necessitates further fine-tuning. To address this, we fine-tune LLaMA2-7B employing our optimized RoPE configuration for a 64k context window (details available in Appendix). As demonstrated in Table~\ref{tbl:finetune}, our approach achieves markedly superior results compared to both PI and YaRN, outperforming these baselines both prior to and following the fine-tuning of LLaMA2-7B. This improvement stems from leveraging non-uniform positional interpolation, which reduces the degree of information loss and supplies a stronger foundation for subsequent fine-tuning.

\textbf{Summary}. This work identifies two kinds of non-uniformity: differences across RoPE dimensions and disparities in token positions. By capitalizing on these non-uniformities during positional interpolation, we achieve substantial enhancements in large language model context extension capability.

\section{LongRoPE}

\label{sec:method}
Building on these observations, we propose LongRoPE, a method that initially implements a robust search algorithm designed to leverage both forms of non-uniformity. Utilizing this approach, we are able to expand the context window of LLMs to exceed 2 million tokens.

\subsection{Problem Formulation}

These dual forms of non-uniformity significantly expand the solution space, thereby increasing the challenges associated with optimization. To tackle this issue, we recast the multidimensional non-uniform position interpolation optimization as a search-based problem.

For a LLM targeting a  context window size of $L'$ and lengthy input documents $\mathbf{X}$, where each $\mathbf{x}\in\mathbf{X}$ surpasses $L'$ in token length, we denote the original rotary angle of the $i^{th}$ dimension in RoPE embedding at token position $n$ as  $\frac{n}{\beta^i}$. The optimization problem is then formulated as follows:

\begin{equation}
\begin{gathered}
		\label{eq:problem}

\underset {\mathbf{x}\in\mathbf{X}; \, |\mathbf{x}|\ge L'}{ \operatorname {arg\,min} } \,  \mathcal{L}\, \left( \text{LLM} (\text{RoPE}, \mathbf{X})\right),	\text{where \;} 
\\
\scalebox{0.93}{
$\underset {\begin{subarray}{l} i=0,\cdot\cdot,\frac{d}{2}-1;\\ n\in[ 0,|\mathbf{x}|);\end{subarray}}{\text{RoPE}_(n)}= {\left[\cdot\cdot,\cos\left(\mathbb{I}(\hat{\lambda}_{i}, \hat{n})\cdot\frac{n}{\beta^i}\right),  \sin\left(\mathbb{I}(\hat{\lambda}_{i}, \hat{n})\cdot\frac{n}{\beta^i}\right),\cdot\cdot\right] }$}\\
	\text{where \;} \mathbb{I}(\hat{\lambda}_{i},\hat{n})=\begin{cases}
	1&\mbox{\;  $n<\hat{n}$}\\
{\frac{1}{\lambda}_{i}}&\mbox{\; $n\geq\hat{n}$}
\end{cases}
	\end{gathered}
\end{equation}

In this approach, we define a set of scaling factors, $\mathbb{I}(\hat{\lambda}_{i},\hat{n})$, to address both types of non-uniformities. Here, $\hat{\lambda_i}$ represents the non-uniformity across RoPE dimensions, and $\hat{n}$ refers to the non-uniformity associated with token positions. In detail, the scaling factor $\mathbb{I}(\hat{\lambda}_{i},\hat{n})$ is applied to modify the rotation angle within the $i^{th}$ RoPE dimension, with $\hat{\lambda}_{i}$ serving as the scaling coefficient and $\hat{n}$ acting as the threshold for token position. For the first $\hat{n}$-1 token positions, the scaling factor $\hat{\lambda}_{i}$ is inactive, and the standard RoPE rotation angle $\frac{n}{\beta^i}$ remains in use. However, once the token position satisfies $n \ge \hat{n}$, the rescaling is employed.

Assuming a desired context window length of $L'$, the task is to determine the most effective rescaling coefficients ($\mathbb{I}(\hat{\lambda}_{0},\hat{n})$, $\mathbb{I}(\hat{\lambda}_{1},\hat{n})$,..., $\mathbb{I}(\hat{\lambda}_{i},\hat{n})$, ...) corresponding to each RoPE dimension from the $1^{st}$ through the $d^{th}$. By applying these optimized rescale factors to the target $\text{LLM}$'s $\text{RoPE}$, our aim is for the model to obtain the lowest possible value of the next token prediction loss, $\mathcal{L}$ (equivalent to perplexity), when processing input sequences $\mathbf{X}$ consisting of $L'$ tokens.

\subsection{Searching the Non-uniform Position Interpolation}

To address the issue presented in Eq.~\ref{eq:problem}, we propose a straightforward and efficient technique that identifies the optimal RoPE rescaling parameters, thereby maximizing the utilization of multidimensional irregularities within the position embedding.

\textbf{Search space}. We establish an extensive search space that accommodates both types of non-uniformities. The composition of this space is summarized in Table~\ref{tbl:searchspace}. More precisely, we permit individualized rescale factors to be selected for every dimension within the RoPE embedding. To simplify the structure of the search space, the search procedure is performed over $\lambda_i$ and $\hat{n}$, rather than directly over $\mathbb{I}(\hat{\lambda}_{i},\hat{n})$, with $\hat{\lambda}_{i}$ defined as $1/\lambda_i$.

As indicated in Table~\ref{tbl:searchspace}, the parameter $\lambda_i$ may assume values ranging from 1.0, corresponding to straightforward extrapolation, up to $s\times1.25$, which reflects an interpolation magnitude exceeding that used in PI. The search increment for $\lambda_i$ is set at 0.01, where $s$ denotes the intended ratio of context window extension.

The parameter $\hat{n}$ specifies how many of the initial token positions preserve the standard RoPE embedding without applying positional interpolation. In practice, we enable $\hat{n}$ to take values from the set \{0, 1, 2, 4, 8, 12, 16, 20, 24, 28, 32, 64, 128, 256\}. Setting $\hat{n}=0$ implies that every token position utilizes the optimized rescale factors.

\textbf{Evolution-based search}. As outlined in Table~\ref{tbl:searchspace}, our search space encompasses a vast array of positional interpolation options, making thorough exploration a complex task. For instance, setting $s=4\times$ results in $400^{128/2}\times14=4\times10^{167}$ possible configurations. The size of the search space grows exponentially as the extension ratio increases. To efficiently navigate this large space, we apply an evolution search strategy~\cite{guo2020single} and incorporate two optimization methods to significantly enhance the efficiency of the search. The overall process is summarized in Algorithm~\ref{alg:search}.

\textit{Optimized initial population generation}. Rather than generating all $P$ rescale factors in the population at random, we explicitly include the three RoPE rescale factors corresponding to PI, NTK, and YaRN as members of the initial group. The rest of the $P-3$ individuals are produced by applying random mutations to these three rescale factors, with each mutation occurring with probability $p$.

\textit{Monotonically non-decreasing constraint}. Once the initial population has been produced, we evaluate the perplexity of each individual by leveraging the corresponding RoPE rescale factors with the target LLM and measuring the perplexity on input $\mathbf{X}$. The top-k scoring individuals are selected to serve as parents in the evolutionary process. Nevertheless, due to the expansive search space, straightforward application of mutation and crossover can result in the exploration of suboptimal candidates, thereby incurring unnecessary perplexity evaluations. This inefficiency is exacerbated in scenarios where $L'$ is large, as computing perplexity requires significant inference time for each individual.

To tackle this issue, we enforce a constraint ensuring that the RoPE rescaling factors sampled are monotonically non-decreasing, such that $\lambda_i\le \lambda_{i+1}$. Only those RoPE variants satisfying this monotonicity condition are utilized during LLM perplexity assessments, which greatly narrows down the search space and enhances efficiency. More precisely, we stipulate that $\lambda_i$ should increase consistently as the RoPE dimension index $i$ progresses from $0$ to $63$. The rationale for enforcing monotonicity across dimensions is anchored in NTK theory~\cite{jacot2018neural,tancik2020fourier,ntk}. According to this theoretical perspective, components at lower dimensions—which represent higher frequencies—require minimal interpolation (reflected by smaller $\lambda_i$), whereas those at higher dimensions, corresponding to lower frequencies, benefit from increased interpolation (thus, larger $\lambda_i$).

\begin{algorithm}[t]
	\small
	\caption{The search algorithm}
	\textbf{Input:} target LLM, input samples $\mathbf{X}$, population size $P$, mutation size $N_1$, crossover size $N_2$, maximum iterations $\mathcal{T}$, mutation probability $p$\\

	\begin{algorithmic}[1]
		\label{alg:search}
		\STATE Top-k $\gets \phi$;
		\STATE $\text{P}_0 \gets$ \textit{Initialize\_population\_with\_optimization} ($P$, $\mathbf{X}$, $p$);
		\FOR{$i = 1$ to $\mathcal{T}$}
		\STATE \textit{Compute\_perplexity} (LLM, $\text{P}_{i-1}$, $\mathbf{X}$);
		\STATE Top-k $\gets$ \textit{Update\_Topk} (Top-k, $\text{P}_{i-1}$);
		\STATE $\text{P}_{mutation} \gets$ \textit{Mutation\_with\_mono\_constraint} (Top-k, $N_1$, $p$);
		\STATE $\text{P}_{crossover} \gets$ \textit{Crossover\_with\_mono\_constraint} (Top-k, $N_2$);
		\STATE $\text{P}_i = \text{P}_{mutation} \cup \text{P}_{crossover} \cup$ Top-k;
		\ENDFOR
		\STATE Output the member of Top-k that achieves the minimum perplexity;
	\end{algorithmic}
\end{algorithm}

\textbf{8$\times$ extension without fine-tuning}. Through evolutionary search, our approach efficiently uncovers non-uniform RoPE rescale factors that selectively maintain critical dimensions and positions, thereby reducing information loss caused by interpolation. As shown in Fig.\ref{fig:non-ft}, this enables our method to expand LLaMA2’s context window from 4k up to 32k without the need for fine-tuning. By comparison, prior methods such as PI and non-uniform NTK or YaRN experience a sharp rise in perplexity when the context is extended beyond 2$\times$.

\subsection{Extending LLM Context Window to 2048K}
\label{sec:extendto2048k}

\textbf{Progressive extension to 2048k}. Here, we present our approach for scaling the context length of pre-trained LLMs from the conventional 4k to more than 2048k tokens. As shown, applying non-uniform positional interpolation enables us to expand the context window by a factor of 8$\times$ without the need for additional fine-tuning. When even greater extensions are needed (such as 512$\times$), fine-tuning becomes essential. One approach involves determining suitable RoPE rescaling factors specific to the desired 2048k context length before proceeding with fine-tuning. However, this strategy is often constrained by the significant computational costs required. In addition, our observations indicate that effectively fine-tuning LLMs for such a substantial extension factor presents notable difficulties (refer to Appendix).

Fortunately, LongRoPE proves to be efficient for application to both the base model and its fine-tuned, longer-context version. Accordingly, we propose a streamlined and incremental approach that reaches the desired 2048k target using only 1k fine-tuning iterations, all conducted at a training context length capped at 256k.

$\Diamond$ \textit{LongRoPE search for extending pre-trained LLMs to 256k}. Using LLaMA2 as a case study, we explore scaling up to context window sizes of 128k and 256k through targeted search. At these configurations, the model is expanded by factors of 32$\times$ and 64$\times$, accordingly.

$\Diamond$ \textit{Fine-tuning to 256k}. Subsequently, we adapt the pre-trained LLM to support a context window of 256k through a staged fine-tuning process. In the initial phase, LLaMA2 undergoes 400 fine-tuning steps utilizing RoPE rescaled factors set for 128k. Once this stage is complete, we update the model by switching to 256k RoPE rescaled factors at the checkpoint, followed by a further 600 steps of fine-tuning. This two-phase approach demonstrates superior efficiency compared to conducting all fine-tuning steps directly at 256k.

$\Diamond$ \textit{Scaling fine-tuned extended LLM to 2048k using LongRoPE search}. In the final step, we apply an additional search phase to the previously fine-tuned LLM with a 256k context length. This enables the model to accommodate an exceptionally large context window of 2048k, achieved without any extra rounds of fine-tuning. Overall, the model's context length is expanded by a factor of $512\times$.

\textbf{Recovery within the original context window}. When the context window is expanded dramatically to 2048k, we observe reduced performance inside the model's initial context range. This phenomenon has been previously reported with positional interpolation~\cite{pi}, which constrains the position embedding space in the original window, forcing it into a comparatively compressed subregion at higher embedding dimensions. Such compression can impair model effectiveness. Notably, when the extension ratio reaches 512$\times$, the position encodings for the original 4k context window are especially densely packed.

To address this issue, we conduct an additional evolution search on the augmented LLM, specifically tuning the RoPE rescale parameters for shorter context lengths such as 4k and 8k. Since shorter sequences require less positional interpolation, we constrain the upper limit of the search space for $\lambda$. At inference time, the LLM adapts the relevant RoPE rescale factors on the fly.

\section{LongRoPE}

\label{sec:method}
Building upon these insights, we propose LongRoPE. This approach initially develops an effective search algorithm that leverages both forms of non-uniformity, subsequently applying it to broaden the context window of LLMs to over 2 million tokens.

\subsection{Problem Formulation}

These two types of non-uniformity can significantly expand the solution space and complicate the optimization process. To tackle this challenge, we recast the optimization of multidimensional non-uniform position interpolation as a search problem.

For a LLM targeting a  context window size of $L'$ and lengthy input documents $\mathbf{X}$, where each $\mathbf{x}\in\mathbf{X}$ surpasses $L'$ in token length, we denote the original rotary angle of the $i^{th}$ dimension in RoPE embedding at token position $n$ as  $\frac{n}{\beta^i}$. The optimization problem is then formulated as follows:

\begin{equation}
\begin{gathered}
		\label{eq:problem}

\underset {\mathbf{x}\in\mathbf{X};\,|\mathbf{x}|\ge L'}{\operatorname{arg\,min}}\,\mathcal{L}\left(\text{LLM}(\text{RoPE},\mathbf{X})\right),\text{ where }\\
\scalebox{0.93}{
$\underset {\begin{subarray}{l} i=0,\dots,\frac{d}{2}-1;\\ n\in [0,|\mathbf{x}|);\end{subarray}}{\text{RoPE}_{(n)}}={[\ldots, \cos(\mathbb{I}(\hat{\lambda}_{i},\hat{n})\cdot \frac{n}{\beta^i}),\ \sin(\mathbb{I}(\hat{\lambda}_{i},\hat{n})\cdot \frac{n}{\beta^i}), \ldots ]}$}
\\
\text{with}~ \mathbb{I}(\hat{\lambda}_{i},\hat{n})=
\begin{cases}
1 &\mbox{ if $n < \hat{n}$}\\
\frac{1}{\lambda_{i}} &\mbox{ if $n \geq \hat{n}$}
\end{cases}
\end{gathered}
\end{equation}

We define a collection of rescale factors, $\mathbb{I}(\hat{\lambda}_{i},\hat{n})$, designed to address both types of non-uniformity present. Here, $\hat{\lambda_i}$ corresponds to non-uniformity across RoPE dimensions, while $\hat{n}$ pertains to positional non-uniformity among tokens. More precisely, $\mathbb{I}(\hat{\lambda}_{i},\hat{n})$ is employed to adjust the rotation angle in the $i^{th}$ RoPE dimension, with $\hat\lambda_{i}$ serving as the scaling parameter and $\hat{n}$ representing the threshold position for tokens. For the first $\hat{n}$-1 token positions, the standard RoPE rotary angle $\frac{n}{\beta^i}$ remains unmodified, meaning the scaling factor $\hat\lambda_{i}$ is not yet applied. Once token positions satisfy $n\ge\hat{n}$, the appropriate rescale factor begins to modify the angle accordingly.

Assuming a desired context window length of $L'$, our aim is to determine the optimal rescaling factors ($\mathbb{I}(\hat{\lambda}_{0},\hat{n})$, $\mathbb{I}(\hat{\lambda}_{1},\hat{n})$, ..., $\mathbb{I}(\hat{\lambda}_{i},\hat{n})$, ...) across the $1^{st}$ through $d^{th}$ dimensions of RoPE. Consequently, applying these rescaled $\text{RoPE}$ parameters to the target $\text{LLM}$ should result in the lowest possible next-token prediction loss $\mathcal{L}$ (perplexity) when processing input sequences $\mathbf{X}$ of length $L'$.

\subsection{Searching the Non-uniform Position Interpolation}

To address the issue outlined in Eq.~\ref{eq:problem}, we propose a straightforward but powerful approach that identifies optimal RoPE rescaling factors, thereby leveraging the complex, multidimensional variations inherent in position embeddings.

\textbf{Search space}. We construct an extensive search space to incorporate both forms of non-uniformity. The search space is detailed in Table~\ref{tbl:searchspace}. In particular, our approach enables the optimization of individual rescale factors for each dimension within the RoPE embedding. To streamline the design of the search space, we choose to directly search over $\lambda_i$ and $\hat{n}$ rather than over $\mathbb{I}(\hat{\lambda}_{i},\hat{n})$, noting that $\hat{\lambda}_{i}=1/\lambda_i$. According to Table~\ref{tbl:searchspace}, the search range for each $\lambda_i$ spans from a lower bound of 1.0 (corresponding to exact extrapolation) up to a maximum of $s\times1.25$ (which denotes a degree of interpolation surpassing that of PI), with increments of 0.01, where $s$ represents the intended extension ratio of the context window.

The parameter $\hat{n}$ specifies how many leading token positions preserve the unaltered positional encoding, specifically utilizing the original RoPE embedding rather than interpolation. In our experiments, $\hat{n}$ is permitted to take values from the set \{0, 1, 2, 4, 8, 12, 16, 20, 24, 28, 32, 64, 128, 256\}. Setting $\hat{n}=0$ indicates that every token position employs the optimized rescale factors.

\textbf{Evolution-based search}. The search space summarized in Table~\ref{tbl:searchspace} covers a vast array of positional interpolation options, making thorough and efficient search a substantial challenge. For instance, when the extension ratio is $s=4\times$, the resulting space encompasses $400^{128/2}\times14=4\times10^{167}$ possible combinations. As the extension ratio increases, the search space grows at an exponential rate. To efficiently traverse this space, we adopt an evolution search strategy as in~\cite{guo2020single} and incorporate two optimization methods to significantly enhance the search process. The overall methodology is described in Algorithm~\ref{alg:search}.

\textit{Optimized initial population generation}. Rather than generating the entire initial population of $P$ rescale factors at random, we explicitly include the three RoPE rescale factors that correspond to PI, NTK, and YaRN as distinct individuals within the initial pool. The rest of the population, totaling $P$-3 individuals, is formed by applying random mutations to these three reference rescale factors, each mutation occurring with probability $p$.

\textit{Monotonically non-decreasing constraint}. Once the initial population has been produced, we evaluate the LLM perplexity of each candidate. For this, we implement the relevant RoPE rescale factors in the target LLM and assess perplexity on input $\mathbf{X}$. The highest-scoring $k$ individuals are selected as parents for the evolutionary process. Nevertheless, due to the expansive search space, unrestrained mutation and crossover can frequently generate low-quality candidates, resulting in excessive and redundant perplexity evaluations. This inefficiency is exacerbated when $L'$ assumes large values, as each perplexity inference requires significant computational effort.

To tackle this challenge, we enforce a non-decreasing monotonicity requirement on the sampled RoPE rescaled factors: $\lambda_i \leq \lambda_{i+1}$. Only those RoPE parameterizations adhering to this constraint are considered for perplexity assessment in LLMs, thus greatly reducing the computational burden of the search space. More concretely, we stipulate that $\lambda_i$ must be a monotonically increasing function of the RoPE dimension ($i = 0,\ldots,63$). This choice of enforcing monotonicity across dimensions draws on insights from NTK theory~\cite{jacot2018neural,tancik2020fourier,ntk}, which implies that lower-indexed (higher frequency) dimensions should undergo minimal interpolation (i.e., have a smaller $\lambda_i$), while higher-indexed (lower frequency) dimensions permit greater interpolation (i.e., a larger $\lambda_i$).

\begin{algorithm}[t]
	\small
	\caption{The search algorithm}
	\textbf{Input:} target LLM, input samples $\mathbf{X}$, population size $P$, mutation size $N_1$, crossover size $N_2$, max iterations $\mathcal{T}$, mutate probability $p$\\

	\begin{algorithmic}[1]
		\label{alg:search}
		\STATE Top-k $\gets \phi$;
		\STATE $\text{P}_0 \gets$ \textit{Initialize\_population\_with\_optimization}$(P, \mathbf{X}, p)$;
		\FOR{$i=1$ to $\mathcal{T}$}
		\STATE \textit{Compute\_perplexity} (LLM, $\text{P}_{i-1}$, $\mathbf{X}$);
		\STATE Top-k $\gets$ \textit{Update\_Topk} (Top-k, $\text{P}_{i-1}$);
		\STATE $\text{P}_{mutation} \gets$ \textit{Mutation\_with\_mono\_constraint} (Top-k, $N_1$, $p$);
		\STATE $\text{P}_{crossover} \gets$ \textit{Crossover\_with\_mono\_constraint} (Top-k, $N_2$);
		\STATE $\text{P}_i \gets \text{P}_{mutation} \cup \text{P}_{crossover} \cup \text{Top-k}$;
		\ENDFOR
		\STATE Output the member in Top-k exhibiting the minimal perplexity;
	\end{algorithmic}
\end{algorithm}

\textbf{8$\times$ extension without fine-tuning}. By leveraging our evolutionary search approach, we successfully determine non-uniform RoPE rescale factors that selectively maintain crucial dimensions and positions, thereby reducing the information degradation caused by interpolation. As shown in Fig.\ref{fig:non-ft}, this technique enables the expansion of LLaMA2's context window from 4k to 32k tokens without necessitating fine-tuning. In comparison, prior approaches such as PI, along with non-uniform NTK and YaRN, experience significant increases in perplexity once the context length is doubled.

\subsection{Extending LLM Context Window to 2048K}
\label{sec:extendto2048k}

\textbf{Progressive extension to 2048k}. In this section, we present our approach for enlarging the context window of pre-trained LLMs from the standard 4k up to more than 2048k. Our non-uniform positional interpolation technique enables an 8$\times$ increase in context length without requiring fine-tuning. To achieve more substantial scaling, such as 512$\times$, fine-tuning becomes essential. A commonly considered strategy involves identifying suitable RoPE rescaling factors for the desired 2048k window and subsequently performing fine-tuning. Nonetheless, this approach is hindered by the high computational and resource costs associated with such extensive training. Additionally, as we have observed, effective fine-tuning at such large extension ratios presents significant difficulties (refer to Appendix).

Fortunately, LongRoPE performs well on both the base and fine-tuned extended LLMs. Consequently, we propose a progressive and efficient strategy that reaches the intended 2048k context window in only 1k fine-tuning iterations, all while remaining within a training length of 256k.

$\Diamond$ \textit{LongRoPE search for expanding pre-trained LLM context to 256k.} Using LLaMA2 as a case study, we perform search operations for context window sizes of both 128k and 256k tokens. At this point, the context expansion ratios achieved are 32$\times$ and 64$\times$ accordingly.

$\Diamond$ \textit{Fine-tuning to 256k}. Subsequently, we adapt the pre-trained LLM to handle a context window of 256k through fine-tuning. Initially, LLaMA2 is fine-tuned for 400 steps using RoPE rescaling factors set for 128k. Following this phase, we swap in the RoPE rescaling factors corresponding to 256k at the resulting checkpoint and carry out a further 600 fine-tuning steps. Compared to directly fine-tuning at 256k, this staged strategy is demonstrated to be more efficient.

$\Diamond$ \textit{Applying LongRoPE search to further expand a fine-tuned extended LLM to 2048k}. In the final step, we conduct an additional search phase utilizing the previously fine-tuned LLM at 256k length. This procedure enables the context window to reach an impressive size of 2048k, achieved without the need for any more fine-tuning. Consequently, the model attains a total extension factor of $512\times$.

\textbf{Recovery within shorter context windows}. Upon increasing the context window to a massive 2048k tokens, we observe that performance deteriorates within the initial window size. This phenomenon has been previously documented for positional interpolation~\cite{pi}, as expanding the context in this way results in the higher-dimensional position embeddings within the initial window being compressed into a much smaller space. This spatial constraint negatively impacts the model's capabilities. The issue is especially pronounced at a 512$\times$ extension ratio, where positions that once occupied the original 4k context window become densely packed.

To address this issue, we conduct an additional evolution-based search on the extended LLM, aiming to refine RoPE rescale factors specifically for shorter context windows (such as 4k and 8k tokens). For these shorter sequences, we constrain the upper bound of the searched $\lambda$ values, since the need for extensive positional interpolation diminishes. In inference, the LLM adaptively modifies the relevant RoPE rescale parameters as appropriate.

 \section{Experiments}

\textbf{Evaluation Tasks and Models}. Our experiments involve the implementation of LongRoPE with both LLaMA2-7B and Mistral-7B. To comprehensively assess model capabilities, we focus on three evaluation dimensions: (1) computing the perplexity of large language models with extended context when processing lengthy documents; (2) conducting the Passkey retrieval task, which tests how effectively a model can extract a specified passkey from a large volume of distractor text; and (3) evaluating performance on established LLM benchmarks, using a context window constrained to 4096 tokens.

\textbf{Fine-tuning}. The LLaMA2 model is fine-tuned using a linear learning rate schedule that starts at 2e-5, utilizing a global batch size of 32. The model is trained for 400 steps on the Redpajama~\cite{redpajama} dataset, where data are divided into segments of 128k tokens, each segment wrapped with the BOS and EOS tokens. Building on the resulting checkpoint, a further 600 steps of fine-tuning are conducted to extend the context window to 256k tokens. Training on the 128k context window employs 8 A100 GPUs via the distributed training framework~\cite{cube}, whereas scaling up to 256k context requires 16 A100 GPUs. For Mistral, a fixed learning rate of 1e-6 and a global batch size of 64 are applied. Fine-tuning for both the 128k and 256k variants adopts the procedure outlined in YaRN~\cite{yarn}, consisting of 400 steps on the Together Computer's Long-Data Collections~\cite{mistraldata} with 16k token sequences. This setup leverages 4 A100 GPUs.

\textbf{Search}. When the target window size is at most 256k, we set $P$=64, $N_1$=$N_2$=16, $p$=0.3, and $\mathcal{T}$=40, and retain the top-32 candidates for mutation and crossover during each cycle. Perplexity is assessed using 5 randomly chosen samples from the PG19 validation set, ensuring that each sample meets the minimum required context length. For window sizes above 512k, both the population and the numbers for mutation and crossover are reduced by half. In this larger setting, perplexity evaluations are conducted on 3 randomly drawn examples from the Pile-Books3~\cite{pile} validation set.

\textbf{Baselines}. To achieve a context window size of 2048k, we conducted fine-tuning on models with 128k and 256k context lengths, resulting in LongRoPE-2048k (ft=128k) for LLaMA2 and LongRoPE-2048k (ft=256k) for Mistral. These two models are evaluated alongside other leading models capable of extending context windows, focusing on open-source LLMs that have undergone further fine-tuning with positional interpolation techniques, including PI, NTK, and YaRN. The comparison covers models such as Together-32k~\cite{together}, Code LLaMA~\cite{codellama}, LongLoRA-full-FT-100k~\cite{longlora}, as well as YaRN-LLaMA and YaRN-Mistral~\cite{yarn}.

\subsection{Main Results}

\label{sec:mainresult}
\textbf{Long sequence language modeling up to 256k tokens}. To start, we evaluate our approach against leading extended LLMs for input sequences with a maximum length of 256k tokens. To showcase the breadth of our method, we draw on two benchmarks: the test partitions of Proof-pile~\cite{pg19} and PG19~\cite{pile}. Our analysis centers on perplexity across a range of context window sizes, employing a sliding window of 256 tokens. For the PG19 dataset, we utilize the full test set, comprising 100 documents. In the case of Proof-pile, we adhere to the protocol set forth by YaRN~\cite{yarn}, selecting at random 10 samples, each no shorter than 128k tokens.

Table~\ref{tbl:proof-pile} and Table~\ref{tbl:pg19} present a comparison of perplexity results for LLaMA2 and Mistral, each augmented with various interpolation strategies, on the Proof-pile and PG19 datasets, respectively. We emphasize two main findings: \textbf{(1)} our extended models exhibit a consistent reduction in perplexity as the evaluation length increases from 4k up to 256k, indicating their effective use of extended context. \textbf{(2)} Notably, even when the models are evaluated at a context window 16$\times$ greater than the original, which generally poses difficulties for preserving performance at more modest lengths, our LongRoPE-2048k variants consistently surpass state-of-the-art baselines within the 256k context range.

\textbf{Long sequence language modeling beyond 2000k}. To assess model performance on exceptionally lengthy texts, we utilize the Books3~\cite{pile} dataset. For efficient evaluation, 20 books are randomly chosen, each with a length surpassing 2048k, and a 256k sliding window is applied.

Table~\ref{tbl:books3} demonstrates that LongRoPE effectively expands the context window of both LLaMA2-7B and Mistral-7B to 2048k, while preserving perplexity levels that are on par with or even better than the baselines across context ranges of 8k to 128k. Distinct performance patterns are evident between LLaMA2 and Mistral at 2048k: Mistral surpasses baseline models for shorter sequences, but its perplexity rises above 7 once contexts exceed 256k. In contrast, LLaMA2 exhibits a predictable trend where perplexity steadily decreases as the context lengthens, with only slight upticks at 1024k and 2048k. Further, LongRoPE-2048k yields improved results on LLaMA2 when fine-tuned with a 256k window compared to 128k, likely because of a smaller extension ratio (8$\times$ versus 16$\times$). Conversely, Mistral achieves better outcomes with a fine-tuning window of 128k. This difference is attributed primarily to our adoption of YaRN’s protocol: for Mistral’s 128k and 256k fine-tuning, a training length of 16k is used, which constrains Mistral’s capacity for further expanding the context window post fine-tuning.

\textbf{Passkey retrieval}. Here, we examine how well the effective context window supports generation tasks. We adopt a synthetic evaluation task for passkey retrieval, as introduced by ~\cite{passkey}. This involves prompting the model to identify a randomly embedded passkey—specifically, a five-digit numeral—concealed within an extensive document. The precise construction of the prompt is provided in the appendix. For our evaluation, we conduct 10 runs of the passkey retrieval task, each time embedding the passkey at a location randomly and uniformly selected throughout the evaluation context window.

Fig.~\ref{fig:passkey} presents a comparison of retrieval accuracy against baseline models. While the retrieval accuracy of prior approaches quickly falls to zero once the context length surpasses 128k, our LongRoPE-LLaMA2-2048k (ft=256k) demonstrates consistently strong performance, sustaining high retrieval accuracy ($\ge$90\%) from 4k up to 2048k, even in the highly demanding setting of extracting a passkey from millions of tokens. For LongRoPE-Mistral-2048k (ft=128k), the model maintains perfect accuracy (100\%) up to 1800k tokens, before decreasing to 60\% at 2048k—an outcome that mirrors the slight rise in perplexity seen at 2048k in Table~\ref{tbl:books3}.

\textbf{Baseline assessments within the native context window}. We assess LongRoPE-2048k models on their default context window employing the Hugging Face Open LLM Leaderboard~\cite{huggingfaceboard}, considering both zero-shot and few-shot conditions. Specifically, our evaluation protocols use 25-shot ARC-Challenge~\cite{allenai:arc}, 10-shot HellaSwag~\cite{zellers2019hellaswag}, 5-shot MMLU~\cite{mmlu}, and 0-shot TruthfulQA~\cite{lin2021truthfulqa}.

Table~\ref{tbl:commonsense} indicates that our models perform similarly to previous models on the standard benchmark, which was initially developed for shorter context lengths, and surpass the original Mistral on TruthfulQA by a margin of +0.5\%. Although LongRoPE-LLaMA2-2048k, which underwent fine-tuning at 256k, exhibits marginally greater declines in performance, its results for the majority of tasks remain generally acceptable.

\subsection{Ablation Results}

\textbf{Effectiveness of the second positional interpolation}. Within the progressive extension framework, our approach leverages a search algorithm to perform a second, non-uniform positional interpolation on large language models (LLMs) that have already undergone fine-tuning and extension. To assess the utility of this method, we conduct empirical evaluations on our fine-tuned LLaMA2-256k model. We then further extend this model to 512k, 1024k, and 2048k positions, applying both PI and YaRN for comparison. Results summarized in Table~\ref{tbl:secondsearch} indicate that our approach of non-uniform positional interpolation maintains perplexity at a stable rate. Conversely, both PI and YaRN experience a rapid rise in perplexity as the context window size increases.

\textbf{Recovery performance at reduced context lengths.} In order to address performance degradation at shorter context sizes, we recalibrate the RoPE factors in LongRoPE-2048k by applying our search procedure. In this process, we impose lower upper limits on the scale factors during the search, thereby reducing interpolation at short context spans such as 4k and 8k. Table~\ref{tbl:shortperformance} provides a comparison of perplexity scores for LongRoPE-LLaMA2-2048k on the Proof-pile at these short context lengths, together with mean LLM benchmark accuracy. The data highlight a marked enhancement in performance when operating with shorter contexts.

\textbf{Analysis on the two forms of non-uniformities}. To disentangle the individual effects of the two types of non-uniformity, we conduct ablation studies. We design two sets of experiments: (i) scaling LLaMA2-7B to shorter sequences of 16k and 32k using various strategies—namely PI, searching solely over RoPE dimension, and searching over both non-uniformities; (ii) expanding our 256k-length fine-tuned LLaMA2 model to 2048k by the same approach. For both, perplexity is assessed without additional fine-tuning. According to Table~\ref{tbl:searchdimension}, introducing non-uniformity in the RoPE dimension markedly lowers perplexity in comparison to PI's linear interpolation. Incorporating non-uniformity in token position offers notable gains at the 16k and 32k lengths, while its effect diminishes at 2048k, perhaps because such long context lengths lessen its relevance. Additionally, retaining only the original tokens (without interpolation) is not effective under these settings, leaving this aspect for exploration in future research.

\section{Related Works}

Apart from position interpolation-based techniques, this section examines alternative strategies found in the literature.

\textbf{Retrieval-based} approaches incorporate an external memory component, storing extended historical context and utilizing retrieval mechanisms to access pertinent documents during inference~\cite{longllama,wang2023augmenting,borgeaud2022improving}. Such approaches usually require significant architectural changes to the underlying LLMs. By comparison, our method is considerably more lightweight, involving only minimal modifications to positional embeddings. Furthermore, it is versatile enough to address a broader range of long-context tasks, including but not limited to long document summarization and few-shot learning, extending beyond the traditional retrieval paradigms.

\textbf{Attention-based context window extensions}. In addition to positional embedding interpolation, certain studies have extended the input context window within the inherent limitations of the original LLM by altering attention mechanisms~\cite{han2023lminfinite, streamingllm,parallelwindow}. Central to this approach is the use of innovative attention masks, which help alleviate the problem of escalating attention costs associated with novel positions. These strategies are complementary to methods involving positional interpolation.

\textbf{Fine-tuning based approaches} investigate methods for optimizing the fine-tuning of pre-trained LLMs equipped with modified position embeddings to handle extended context windows. Research such as Code LLaMA~\cite{codellama}, LLaMA2 Long~\cite{xiong2023effective}, and ScaledRoPE~\cite{liu2023scaling} typically opts for an especially large base value within RoPE, followed by fine-tuning the model on the desired longer context lengths. In contrast, our technique supports a wide range of target lengths, even surpassing the 2M mark. Furthermore, with context lengths exceeding 128k requiring significant GPU capacities for training, recent approaches like LongLoRA~\cite{longlora} and PoSE~\cite{zhu2023pose} have emerged to alleviate such resource constraints. Our proposed method is complementary to these advances in efficient fine-tuning.

\section{Conclusion}

In this study, we introduce LongRoPE, a technique that significantly extends the context window of LLMs to a record-setting 2048k tokens, all while preserving their original capabilities for shorter contexts. Our approach leverages two distinct types of non-uniformity in RoPE positional embeddings through an efficient evolutionary search process. This dual-purpose strategy both supplies robust initialization for subsequent fine-tuning and allows for an 8$\times$ enlargement of the context window even in the absence of fine-tuning. Furthermore, we outline a progressive extension framework that employs 256k-token fine-tuned LLMs to incrementally achieve a 2048k-token context length, requiring no additional fine-tuning. Comprehensive experimental results confirm the efficacy of LongRoPE. We anticipate that our LongRoPE-2048k models will facilitate a wide range of novel long-context applications and stimulate further exploration in this field.

\section*{Broader Impacts} The research described in this paper aims to contribute to progress in the area of Machine Learning. Although our work may have various possible effects on society, we have not identified any particular issues that warrant explicit mention in this context.

	\balance

\end{document}